\uuid{Kmr3}
\titre{Écriture algébrique}
\auteur{bodin ; hanani ; mzari}
\organisation{exo7}
\licence{CC BY-NC-SA 4.0}
\datecreate{2020-01-20}
\niveau{L1}
\module{Algèbre}
\chapitre{Nombres complexes}
\sousChapitre{Forme cartésienne, forme polaire}
\difficulte{1}
\theme{Écritures algébrique et géométrique, écriture algébrique}
\exo7qcmid{lille-1.nombres-complexes.002}

\contenu{
\begin{question}
\qclassification{104.01}
\qtags{motcle=écriture algébrique}
Soit $z=\frac{i+1}{1-i\sqrt 3}$. Quelles sont les assertions vraies ?
\begin{answers}
    \good{$|z|=\frac{1}{\sqrt 2}$}

    \good{ $z\overline{z} =\frac{1}{2}$}

    \good{Un argument de $z$ est : $\frac{7\pi}{12}$.}

    \bad{Le conjugué de $z$ est : $\overline{z}=\frac{i-1}{1+i\sqrt 3}$.}

\end{answers}
\begin{explanations}
On applique les formules :
$|\frac{z_1}{z_2}|= \frac{|z_1|}{|z_2|}$, $|z|^2=z\overline{z}$ et $\arg(\frac{z_1}{z_2})= \arg z_1 - \arg z_2 \, [2\pi]$.
\end{explanations}
\end{question}
}
